\documentclass[12pt]{article}

\usepackage[a4paper, total={6in, 8in}]{geometry}
\usepackage{textcomp}

\begin{document}
\setlength{\parindent}{0pt}

\def \t {\textrightarrow}
\def \v {\vspace{1em}}
\def \bi {\begin{itemize}}
\def \ei {\end{itemize}}
\def \s[#1] {\section*{#1}}
\def \ss[#1] {\subsection*{#1}}
\def \sss[#1] {\subsubsection*{#1}}

\s[Il gelsomino notturno]
L'ape rappresenta di per se l'operosita \\
La voce verbale sussurra \t evoca una onomatopea evidente \t sussurro è dell'essere umano e non dell'ape, quindi bisogna anche aggiungere sinestesia \\
Ape è tardiva \t il processo naturale richiede che alcune siano operative di giorno e altre di notte \\
Le celle sono le celle dell'alveare \\
La chioccetta con la maiuscola \t due aspetti: chioccia è protezione, è madre e si vede come evocazione della famiglia \\
Aia azzura \t si confonde con il cielo \\
Pigolio di stelle \t è una sinestesia e onomatopea \t associa alle stelle ciò che è proprio della chioccetta \\
Le stelle sono presenti \t sono spettatrici \\
Nella prossima strofa si passa da sfere visive a quelle olfattive \t l'odore che porta il vento \t che è il profumo del gelsomino, che di notte profuma di piu \\
Il fuori ci accompagna verso il dentro \t a meta della penultima strofa \t il fuori è l'evento, ovvero la festa del matrimonio e la prima notte \\
Si cova / enj. \\
Urna molle e segreta \t si riferisce all'utero \t è il calice della vita \t qua invece viene chiamato urna, che è di morte \\
Il sentore di morte è presente e non abbandona mai l'anima dell'autore \\
Non so che è il mistero della vita, il miracolo della vita \t la nascita è sempre una sorpresa, in quanto immaginare \\
Si legge anche una morbosita sessuale e una morbosita nell utilizzo dei termini che descrivono l'apparato genitale femminile \\
La morbosita si comprende anche nella negazione pascoliana dell'eros \t Recalcati ricorda anche Lacan, per cui l'io (che nel linguaggio lacaniano è una diga) si scontra con le pulsioni dell'eros \t che urtano contro la diga, che vengono rimbalzate \\
Le due pulsioni, eros e tanatos (Freud) portano a una congestione (eccesso) dell'eros, tale che l'unica valvola di sfogo è tanatos (la morte) \\
Si distingue quindi la costrizione o una scelta della castita \t la critica psicanalitica indiviuda un blocco infantile dell'autore, e la regressione verso la madre che si accompagna al rifiuto di un impengo stabile dell'autore \t che significherebbe avere un nido costante

\v

In un altro testo delle Myricae: sconcertante sfogo rispetto all'incapacita di affrontare la vita \t e di rimanere bambino \t si augura una morte pur di vedere in vita la madre = morbosita verso la madre e solitudine non risolvibile

\ss[Domande possibili]
Fogazzaro è inteprete del post-manzonismo \t vale quello che si dira per d'annuznio: per supeare qualcosa bisogna attraversarlo \\
La perosnalita di manzoni è ingombrante, e condiziona la letteratura molto \t la corrente del modernismo si realizza come fattibile quando la scienza avanza, e Compte mostra fiducia nella scienza tradita dai risultati \t uomo non puo essere studiato nell anima \\
Fogazzaro partecipa ai valori del decadentismo \t e mostra la fase di passaggio per cui la conflittualita tra fede e ragione è accessa \\
Quello che si intende oggi come mondo della etica religiosa e della scienza \t la cultura scientifica presente ostruzionismo religioso \\
Alcuni si devono sottoporre a una fecondazone assistita \t chiesa è recalcitante \\
Bio-etica \t le nascite vengono pilotate, e intervento su un patrimonio genetico \t e inoltre la prassi dell'utero in affitto \t rappresenta la conflittualita della scienza e chiesa = modernismo in Fogazzaro, che vive questa tensione \\
Lui è attaccato ai valori, ma la chiesa mostra piegatura verso valori \t modernismo si focalizza su un ansia di rinnovamento, che chiede alla chiesa di ridimensionarsi davanti alla scienza \t si vede per esempio in piccolo mondo moderno 

\v

Come si giustifica la ribellione degli scapigliati? \t il postmanzonismo contesta i paramentri morali di manzoni, fortemente religiosi \t infatti manzoni padroneggia nel romanzo, ma anche come testimone di una sicurezza non piu presente \t unita d italia ha distrutto ogni certezza \t tutti la volevano, ma poi non si realizza come ci si aspettava \\
Questa mancata realizzazione durata per molto con occupazione straniera, ha portato a delusione profonda \t e sfiducia nel manzonismo \\
Noi siamo figli di padri ammalati \t che sono quelli del manzonismo \t parlano di un linguaggio che parla alle coscienze, che pero non trovano soddisfazoine perche unita d italia porta ad unione di realta molto diverse \t ed è un imposizione e non un progetto condiviso

\v

Verga ha sua parentesi scapigliata \t intensa, che sperimenta a Milano \t stravolge il suo orizzonte limitato \t attaccamento alla sua terra, che se tradito porterebbe alla rovina \\
La scapigliatura in "Nedda" e "Rosso malpelo" vede il passaggio da periodo scapigliato alla concretizzazoine di una nuova fase \\
Diventa quindi seme di una progettualita + ampia \t ovvero il ciclo dei vinti \\
Perimetro limitato, nel verismo c'è declassamento dell autore nei personaggi (criterio dell impersonalita) secondo Sciascia \\
Ma è possibile il romanzo naturale in toto \t no, perche l autore sente i personaggi come le sue vite \\
Questo ridefinisce il paradigma verghiano \t siamo attaccati alla terra come l'ostrica allo scoglio \t come i personaggi si staccano, trovano la morte e questo è presente nel primo romanzo con Lia \t si stacca dal suo paese, va in francia, spera di trovare altro (il fuori è la possibilita di riscatto) ma trova la morte dell anima \t deve prostituirsi \\
Porta a un esclusione + ampia e a vedere il percorso della donna, che è anche carme, ombra (Fosca), non è solo angelo ma anche diavolo \t e Lia è macchiata di peccato, che impersonifica questo \\
Bisogna motivare la scelta della lirica, analisi deve essere contenutistica nel lessico \t soprattuto temi, contenuti, relazione con altre liriche
\end{document}
